\section{Introduction}\label{sec:introduction84}
A measurement with a classical output may retain a quantum advantage
when its input is correlated with a reference.  The question considered
here is quantitative: how well can two ordered measurements be
separated with at most $N$ uses, and which features of their effects
determine the dependence on $N$?  The answer must accommodate both
strictly positive effects and singular measurements.  Positivity of
each effect alone is not the relevant constraint: the effects sum to
the identity, and their boundary geometry is coupled by this equation.

\subsection{The model and the principal results}
An ordered $k$-outcome measurement on $\C^d$ is a tuple
$E_j\succeq0$ with $\sum_jE_j=I_d$.  One call is the channel
\[
 \mathcal M_E(\rho)=\sum_{j=1}^k\tr(E_j\rho)|j\rangle\langle j|.
\]
The input is consumed and the output is only the classical label;
no residual device system is available.  A tester may prepare inputs
entangled with its own reference, retain quantum memory, and choose
subsequent controls using previous labels.  It uses the same control
rule under both hypotheses.  Its public stopping time is bounded
by $N$.  We define $D_N(E,F)$ to be the supremum of the unhalved
trace distance between its two final states.  Thus $0\le D_N\le2$.
Nonadaptive testers may have entangled inputs and retained references;
``nonadaptive'' does not mean ``product input.''

Our complete-body upper bound uses a positive, concave covariance
$\mathcal C_E$ on zero-sum Hermitian tuples.  Its factorization through
a horizontal projection gives a finite-use bound valid through changes
of rank and zero effects.  With $M=(E+F)/2$ and $H=F-E$,
\begin{equation}\label{eq:frontupper84}
 D_N(E,F)\le\min\left\{2,2\sqrt{k+1}
  \sqrt{N\langle H,(\mathcal C_M+N^{-1}\operatorname{Id})^{-1}H\rangle}
                              \right\}.
\end{equation}
This is a global upper certificate, not a claimed formula for the
optimal adaptive distance.  The normalization includes the exact
binary Sylvester restriction of Theorem~\ref{thm:closedcovariance83}.

The new support description determines its entire kernel.  If
$P_j=\operatorname{supp}E_j$ and
$\mathcal W_E=\sum_jP_j\mathcal H_dP_j$, then
Theorem~\ref{thm:supportkernel84} parametrizes all null directions
and proves
\[
 \dim\ker(\mathcal C_E|_{\mathcal T})
 =d^2-\dim\mathcal W_E+\sum_j(d-\operatorname{rank}E_j)^2.
\]
This distinction between support span and the mere effect span is
essential for effects of rank greater than one.  In particular,
a nonprojective measurement can have a nontrivial covariance kernel;
the earlier projective characterization concerns a zero operator,
not one null direction.

This kernel has a finite-use interpretation on every fixed unitary
orbit $E_j(t)=e^{itG}E_je^{-itG}$.  Apart from the stationary case,
Theorem~\ref{thm:orbittrichotomy84} proves, uniformly in every integer
$N\ge1$ and all sufficiently small $|t|$,
\begin{equation}\label{eq:frontorbit84}
 D_N(E,E(t))\asymp_{E,G}
 \begin{cases}
 \min\{1,\sqrt N|t|\},&G\in\mathcal W_E,\\
 \min\{1,N|t|\},&G\notin\mathcal W_E.
 \end{cases}
\end{equation}
The second lower bound uses a corrected logical qubit and the explicit
finite-angle estimate \eqref{eq:finiteanglelower84}.
The support span is precisely the Hermitian Kraus-product span of
this measurement channel.  The metrological criterion is credited
to Zhou--Jiang \cite{ZJ84}; the present formulation identifies it
with the covariance kernel and proves the finite-pair trace estimates
directly.  This is an orbit theorem on the complete body, not a
full-boundary volume or minimax theorem.

We also replace a coordinate-fixed ridge by a public scalar-measurement
covariance.  Transporting its probability vector through the same
output channel gives a regularized data-processing inequality and
exact invariance under reversible label splitting
(Theorem~\ref{thm:transportedridge84}).  The original fixed Euclidean
ridge is not asserted to have this property.

\subsection{Learning and descriptions on balanced interiors}
The operational consequences are complemented by exact description
and common learning laws on the balanced interior
$I_d/(2k)\preceq E_j\preceq3I_d/(2k)$.  Its real affine dimension
is $p=(k-1)d^2$.  Theorems~\ref{thm:multioutcomeentropy82} and
\ref{thm:multioutcomecode82} give
\[
 \log_2\mathcal C_{N,k}(\delta)
 =\frac p2\log_2N+p\log_2(1/\delta)+O(p\log(k+1)),
 \qquad 0<\delta<1/128.
\]
The common learner has the joint order
\[
 M^*_{d,k}(N,\delta,\eta)
 =\Theta_k\!\left(N\delta^{-2}[d^2+d\log(1/\eta)]\right)
\]
for fixed $k$, $d,N\ge1$, $0<\delta\le2^{-24}$ and
$0<\eta\le1/8$.  The displayed constructive upper has a $k^3$
factor; the lower does not match growing $k$.
Theorem~\ref{thm:affinerepair83} supplies polynomial, search-free
statistical legalization, while the entropy-optimal public dictionary
remains a separate finite, potentially exhaustive construction.
The lower theorem permits arbitrary coherent adaptive training,
whereas the upper uses fresh acquisitions and collective processing
of completed outputs.  Pair-dependent discrimination witnesses do
not serve as a common learning design.

\subsection{Resource conventions and reading order}
Unknown-device calls, future-use horizon, transmitted index length,
public dictionary reconstruction, classical arithmetic, trusted-control
synthesis, and physical execution are distinct resources.  A query
bound here holds on every record, with a defined legal fallback.
Public dictionaries and their parameters are independent of the
target; target-dependent private advice would count as transmitted
information.  Exact rational certificate evaluation and balanced
legalization have polynomial bit complexity.  Neither statement
asserts polynomial synthesis of an arbitrary collective readout or
of the whole entropy-optimal dictionary.  The orbit recovery is an
ideal-control comparison construction, not an executed common learner.

Sections~\ref{sec:covariance83}--\ref{sec:ridge84} develop the global
upper, support kernel, orbit classification and transported ridge.
Sections~\ref{sec:finiteoutcome82}--\ref{sec:affinerepair83} give the
balanced-interior learning and coding consequences.
All binary auxiliary proofs used by this article are supplied in the
separate \emph{Binary Foundations and Auxiliary Proofs} supplement.
Cross-document theorem references carry the prefix S and are resolved
from its current source, not from a historical PDF.  The primary article
and supplement together are a self-contained proof package; the
independent structural paper is not a premise of any measurement theorem.
